\documentclass[10pt,a4paper]{amsart}
\usepackage{amsmath,amssymb,amsthm,graphicx}
\usepackage{verbatim}
\newenvironment{reference}{\comment}{\endcomment}
%\newenvironment{reference}{}{}
\newenvironment{slogan}{\comment}{\endcomment}
\newenvironment{history}{\comment}{\endcomment}

% For commutative diagrams we use Xy-pic
\usepackage[all]{xy}

% We use 2cell for 2-commutative diagrams.
\xyoption{2cell}
% Theorem environments.
%
\theoremstyle{plain}
\newtheorem{theorem}[subsection]{Theorem}
\newtheorem{proposition}[subsection]{Proposition}
\newtheorem{lemma}[subsection]{Lemma}

\theoremstyle{definition}
\newtheorem{definition}[subsection]{Definition}
\newtheorem{example}[subsection]{Example}
\newtheorem{exercise}[subsection]{Exercise}
\newtheorem{situation}[subsection]{Situation}

\theoremstyle{remark}
\newtheorem{remark}[subsection]{Remark}
\newtheorem{remarks}[subsection]{Remarks}

\numberwithin{equation}{subsection}

% Macros
%
\def\lim{\mathop{\mathrm{lim}}\nolimits}
\def\colim{\mathop{\mathrm{colim}}\nolimits}
\def\Spec{\mathop{\mathrm{Spec}}}
\def\Hom{\mathop{\mathrm{Hom}}\nolimits}
\def\Ext{\mathop{\mathrm{Ext}}\nolimits}
\def\SheafHom{\mathop{\mathcal{H}\!\mathit{om}}\nolimits}
\def\SheafExt{\mathop{\mathcal{E}\!\mathit{xt}}\nolimits}
\def\Sch{\mathit{Sch}}
\def\Mor{\mathop{\mathrm{Mor}}\nolimits}
\def\Ob{\mathop{\mathrm{Ob}}\nolimits}
\def\Sh{\mathop{\mathit{Sh}}\nolimits}
\def\NL{\mathop{N\!L}\nolimits}
\def\CH{\mathop{\mathrm{CH}}\nolimits}
\def\proetale{{pro\text{-}\acute{e}tale}}
\def\etale{{\acute{e}tale}}
\def\QCoh{\mathit{QCoh}}
\def\Ker{\mathop{\mathrm{Ker}}}
\def\Im{\mathop{\mathrm{Im}}}
\def\Coker{\mathop{\mathrm{Coker}}}
\def\Coim{\mathop{\mathrm{Coim}}}

% Boxtimes
%
\DeclareMathSymbol{\boxtimes}{\mathbin}{AMSa}{"02}

%
% Macros for moduli stacks/spaces
%
\def\QCohstack{\mathcal{QC}\!\mathit{oh}}
\def\Cohstack{\mathcal{C}\!\mathit{oh}}
\def\Spacesstack{\mathcal{S}\!\mathit{paces}}
\def\Quotfunctor{\mathrm{Quot}}
\def\Hilbfunctor{\mathrm{Hilb}}
\def\Curvesstack{\mathcal{C}\!\mathit{urves}}
\def\Polarizedstack{\mathcal{P}\!\mathit{olarized}}
\def\Complexesstack{\mathcal{C}\!\mathit{omplexes}}
% \Pic is the operator that assigns to X its picard group, usage \Pic(X)
% \Picardstack_{X/B} denotes the Picard stack of X over B
% \Picardfunctor_{X/B} denotes the Picard functor of X over B
\def\Pic{\mathop{\mathrm{Pic}}\nolimits}
\def\Picardstack{\mathcal{P}\!\mathit{ic}}
\def\Picardfunctor{\mathrm{Pic}}
\def\Deformationcategory{\mathcal{D}\!\mathit{ef}}
\usepackage{cleveref}
\crefname{originallemma}{original lemma}{original lemmas}
\Crefname{originallemma}{Original lemma}{Original lemmas}
\makeatletter
\newlabel{lemma-quasi-compact}{{00E8}{source}{}{sourcepointer-1}{}}
\newlabel{lemma-quasi-compact@cref}{{[originallemma][0][]00E8}{[1][0][]source}{}{}{}}
\makeatother
\crefname{originallemma}{original lemma}{original lemmas}
\Crefname{originallemma}{Original lemma}{Original lemmas}
\makeatletter
\newlabel{lemma-spec-spectral}{{090M}{source}{}{sourcepointer-2}{}}
\newlabel{lemma-spec-spectral@cref}{{[originallemma][0][]090M}{[1][0][]source}{}{}{}}
\makeatother
\crefname{originallemma}{original lemma}{original lemmas}
\Crefname{originallemma}{Original lemma}{Original lemmas}
\makeatletter
\newlabel{lemma-standard-open-containing-maximal-point}{{00EV}{source}{}{sourcepointer-3}{}}
\newlabel{lemma-standard-open-containing-maximal-point@cref}{{[originallemma][0][]00EV}{[1][0][]source}{}{}{}}
\makeatother
\crefname{originallemma}{original lemma}{original lemmas}
\Crefname{originallemma}{Original lemma}{Original lemmas}
\makeatletter
\newlabel{topology-lemma-characterize-profinite-spectral}{{0905}{source}{}{sourcepointer-4}{}}
\newlabel{topology-lemma-characterize-profinite-spectral@cref}{{[originallemma][0][]0905}{[1][0][]source}{}{}{}}
\makeatother
\crefname{originallemma}{original lemma}{original lemmas}
\Crefname{originallemma}{Original lemma}{Original lemmas}
\makeatletter
\newlabel{topology-lemma-profinite}{{08ZY}{source}{}{sourcepointer-5}{}}
\newlabel{topology-lemma-profinite@cref}{{[originallemma][0][]08ZY}{[1][0][]source}{}{}{}}
\makeatother
\title{Nine equivalent conditions for a spectrum to be profinite, Hausdorff or totally disconnected}
\author{The Stacks Project Authors}
\date{stacks-project:algebra.tex:local-snapshot-20261012; GNU FDL 1.2 or later}
\begin{document}
\maketitle
\noindent\textit{Editorial note.} All delivered mathematical text is the author's own, copied from the pinned Stacks Project chapter file algebra.tex; no statement or proof step was rewritten, reordered or removed.\par\medskip
\noindent\textit{Editorial note.} Preamble provenance, as required for this source: the frozen snapshot stages the chapter's own preamble file beside it (preamble.tex, 217 lines, 8116 bytes, the same shared Stacks-wide preamble the published chow and etale-cohomology cards use); the first line of the pinned chapter file is \textbackslash{}input\{preamble\}, so no preamble text lives inside the chapter itself. The card takes that file's lines 8-18 (the verbatim package and the three \textbackslash{}comment-based environments reference/slogan/history at lines 9-12, section 13.2, together with the Xy-pic setup) and its lines 154 to the runtime end of the file (the native theorem declarations and mathematical macros). Not one character was changed, no package was added, no macro was defined, and the excerpt serves only as preamble, never as a span coordinate.\par\medskip
\noindent\textit{Editorial note.} The same-chapter pointers are located at their \textbackslash{}label lines inside the pinned chapter file; the cross-chapter pointers are located in the project's own tag list staged in the snapshot (sources/stacks-chow/tags.txt), at the verbatim lines carrying their canonical tags.\par\medskip
\noindent\textit{Editorial note.} The counted proof refers to same-chapter and Topology chapter statements by label; they are recorded as navigation stubs pointed at their \textbackslash{}label lines, and the arguments using them stay inside the counted span.\par\medskip
\noindent\textit{Editorial note.} The counted spans contain no \textbackslash{}cite, so the card carries no bibliography.\par\medskip
\noindent\textit{Allowed standard prerequisites (editorial).}\begin{itemize}
\item Standard opens and quasi-compact opens of a spectrum. Used at the start of the counted proof (lines 4993-4995). Identifies quasi-compact opens with finite unions of standard opens.
\item Closedness of a standard open containing a maximal point. Invoked in the counted proof at line 4996. Gives one of the implications between the topological conditions.
\end{itemize}\medskip
\ifdefined\hypertarget\hypertarget{sourcepointer-1}{}\fi
\noindent\textit{Source reference.} Original lemma 00E8: sources/stacks-algebra/pinned-original-source.tex, source lines 3469--3469. Navigation pointer for \textbackslash{}ref\{lemma-quasi-compact\} used inside the counted span (line 5012); the argument that consumes it is retained in the span.\par\medskip
\ifdefined\hypertarget\hypertarget{sourcepointer-2}{}\fi
\noindent\textit{Source reference.} Original lemma 090M: sources/stacks-algebra/pinned-original-source.tex, source lines 4924--4924. Navigation pointer for \textbackslash{}ref\{lemma-spec-spectral\} used inside the counted span (line 5017); the argument that consumes it is retained in the span.\par\medskip
\ifdefined\hypertarget\hypertarget{sourcepointer-3}{}\fi
\noindent\textit{Source reference.} Original lemma 00EV: sources/stacks-algebra/pinned-original-source.tex, source lines 4956--4956. Navigation pointer for \textbackslash{}ref\{lemma-standard-open-containing-maximal-point\} used inside the counted span (line 4997); the argument that consumes it is retained in the span.\par\medskip
\ifdefined\hypertarget\hypertarget{sourcepointer-4}{}\fi
\noindent\textit{Source reference.} Original lemma 0905: sources/stacks-chow/tags.txt, source lines 11398--11398. Navigation pointer for \textbackslash{}ref\{topology-lemma-characterize-profinite-spectral\} used inside the counted span (line 5018); the argument that consumes it is retained in the span.\par\medskip
\ifdefined\hypertarget\hypertarget{sourcepointer-5}{}\fi
\noindent\textit{Source reference.} Original lemma 08ZY: sources/stacks-chow/tags.txt, source lines 11391--11391. Navigation pointer for \textbackslash{}ref\{topology-lemma-profinite\} used inside the counted span (line 5013); the argument that consumes it is retained in the span.\par\medskip
\begin{lemma}
\label{lemma-ring-with-only-minimal-primes}
Let $R$ be a ring. Let $X = \Spec(R)$ as a topological space.
The following are equivalent
\begin{enumerate}
\item $X$ is profinite,
\item $X$ is Hausdorff,
\item $X$ is totally disconnected.
\item every quasi-compact open of $X$ is closed,
\item there are no nontrivial inclusions between its prime ideals,
\item every prime ideal is a maximal ideal,
\item every prime ideal is minimal,
\item every standard open $D(f) \subset X$ is closed, and
\item add more here.
\end{enumerate}
\end{lemma}

\begin{proof}
First proof. It is clear that (5), (6), and (7) are equivalent.
It is clear that (4) and (8) are equivalent as every quasi-compact
open is a finite union of standard opens.
The implication (7) $\Rightarrow$ (4) follows from
Lemma \ref{lemma-standard-open-containing-maximal-point}.
Assume (4) holds. Let $\mathfrak p, \mathfrak p'$ be distinct
primes of $R$. Choose an $f \in \mathfrak p'$, $f \not \in \mathfrak p$
(if needed switch $\mathfrak p$ with $\mathfrak p'$).
Then $\mathfrak p' \not \in D(f)$ and $\mathfrak p \in D(f)$.
By (4) the open $D(f)$ is also closed.
Hence $\mathfrak p$ and $\mathfrak p'$ are in disjoint open
neighbourhoods whose union is $X$. Thus $X$ is Hausdorff and totally
disconnected. Thus (4) $\Rightarrow$ (2) and (3).
If (3) holds then there cannot be any specializations
between points of $\Spec(R)$ and we see that (5) holds.
If $X$ is Hausdorff then every point is closed, so (2) implies (6).
Thus (2), (3), (4), (5), (6), (7) and (8) are equivalent.
Any profinite space is Hausdorff, so (1) implies (2).
If $X$ satisfies (2) and (3), then $X$ (being quasi-compact by
Lemma \ref{lemma-quasi-compact}) is profinite by
Topology, Lemma \ref{topology-lemma-profinite}.

\medskip\noindent
Second proof. Besides the equivalence of (4) and (8) this follows
from Lemma \ref{lemma-spec-spectral} and purely topological facts, see
Topology, Lemma \ref{topology-lemma-characterize-profinite-spectral}.
\end{proof}

\end{document}
